\documentclass[FM]{tulpresentation}
\usepackage{tabularx}
\usepackage{booktabs}
\usepackage{xcolor}
\title{Architektura systému pohledové inspekce kvality}
\subtitle{Kevin Daněk}
\author{Vedoucí: Ing. Igor Kopetschke}

\definecolor{Success}{HTML}{008236}
\definecolor{Warning}{HTML}{d08700}
\definecolor{Error}{HTML}{e7000b}

\begin{document}

\TULtitleframe

\begin{frame}
	\frametitle{Zadání}
	\begin{items}
		\item Analyzujte architekturu softwarového řešení ve společnosti CLOUDCODE pro kamerovou inspekci kvality.
		\item Identifikujte a odůvodněte její nedostatky a navrhněte potřebné změny reflektující moderní požadavky softwarového inženýrství.
		\item Implementujte sledování řešení pomocí OpenTelemetry a softwarových řešení společnosti Grafana Labs.
		\item Implementujte automatické testování a sestavení pomocí nástrojů pro CI/CD.
	\end{items}
\end{frame}

\begin{frame}
	\frametitle{Teoretický základ}
	\begin{items}
		\item Technický dluh = Jev, který způsobuje, že údržba systému vyžaduje stále větší úsilí.
		\begin{subitems}
		\item Vzniká tam, kde chybí systematický návrh a ucelená vize.
		\item Je úhrnem všech nešvárů v kódu - velké třídy, dlouhé metody, mrtvý kód, ...
		\item Detekujeme ho pomocí statické analýzy a podle dynamiky vývoje.
		\item Přecházíme mu pomocí uplatnění správných návrhových principů (Clean Code, Návrhové vzory, Architektonické vzory, Integrační vzory)
	 	\end{subitems}
	\end{items}
\end{frame}

\begin{frame}
	\frametitle{Dosažené výsledky}
	\begin{items}
		\item Řešení společnosti CLOUCODE bylo rozděleno na mikroslužby.
		\item Jednotlivé mikroslužby používají patřičné návrhové vzory (kde jsou potřeba):
		\begin{subitems}
			\item \textbf{REST API} - FastAPI - \textit{Builder, Strategy, Observer}
			\item \textbf{Frontend} - React - \textit{Factory, Composite}
			\item \textbf{Peripherals API} - Celery - \textit{Singleton, Strategy, Adapter}
			\item \textbf{Machine Learning API} - Celery - \textit{Builder, Strategy, Adapter}
		\end{subitems}
	\end{items}
\end{frame}

\begin{frame}
	\frametitle{Dosažené výsledky}
	\begin{items}
		\item Komunikace mezi kontejnery napřímo pomocí Docker DNS.
		\item Vysoká propustnost zajištěna pomocí Redis fronty.
		\begin{subitems}
			\item Snímání snímků z kamery - běžně až 5 za vteřinu
			\item Pokud má stanice 4 kamery, zátěž rychle roste.
			\item Je nežádoucí přehltit kontejner s REST API.
		\end{subitems}
		\item Celé řešení je monitorováno a vizualizováno v Grafana Dashboardu.
	\end{items}
\end{frame}

\begin{frame}
	\frametitle{Shrnutí}
	\begin{table}[H]
		\centering
		\renewcommand{\arraystretch}{1.5}
		\begin{tabularx}{\textwidth}{|X|l|}
			\hline
			\textbf{Bod zadání} & \textbf{Stav}\\\hline
			Analýza architektury stávajícího softwarového řešení & \textbf{\textcolor{Success}{Ano}} \\\hline
			Identifikace a odůvodnění změn a návrh řešení & \textbf{\textcolor{Success}{Ano}} \\\hline
			Implementace sledování pomocí OpenTelemetry & \textbf{\textcolor{Success}{Ano}} \\\hline
			Implementace automatického testování a sestavení & \textbf{\textcolor{Warning}{Částečně}} \\\hline 
		\end{tabularx}
	\end{table}
\end{frame}
\TULendframe
\end{document}
